%! Author = matteomagnini
%! Date = 18/03/26

% Preamble
\documentclass[11pt]{article}

% Packages
\usepackage{amsmath}
\usepackage{amssymb}
\usepackage{my-acronyms}
\usepackage{hyperref}

\title{
    Intelligent System 2\\
    Tutorial week 5: Beliefs revision and update
}

\author{
    Matteo Magnini\\
    \href{mailto:matteo.magnini@uni.lu}{matteo.magnini@uni.lu}
}

\date{
    19 March, 2026
}

% Document
\begin{document}

    \maketitle

    \section{Solutions}\label{sec:solutions}

    \subsection{Solution to Tweety}

    Let the propositional variables be:
    \[
    b = \text{Tweety is a bird}, \quad
    p = \text{Tweety is a penguin}, \quad
    f = \text{Tweety flies}
    \]

    The initial belief state is:
    \[
    T = Cn(\{b \rightarrow f,\; p \rightarrow b,\; p\})
    \]

    From $T$ we derive:
    \[
    T \vdash b \quad \text{and} \quad T \vdash f
    \]

    We now learn:
    \[
    A = \neg f
    \]

    Since:
    \[
    T \vdash f \quad \Rightarrow \quad T \vdash \neg A
    \]
    this is a \textbf{non-trivial belief revision}.

    \paragraph{Revision step (AGM intuition).}
    We must remove (at least) one formula responsible for deriving $f$.

    The conflict comes from:
    \[
    b \rightarrow f
    \]

    We keep:
    \[
    p \rightarrow b, \quad p
    \]

    Thus, a maximal consistent subset not implying $f$ is:
    \[
    T' = Cn(\{p \rightarrow b,\; p\})
    \]

    Now expand with $\neg f$:
    \[
    T \star \neg f = Cn(\{p \rightarrow b,\; p,\; \neg f\})
    \]

    \paragraph{Result.}
    \[
    T \star \neg f \vdash p \wedge b \wedge \neg f
    \]

    \textbf{Interpretation:}
    Tweety is still a penguin and a bird, but we no longer believe that birds always fly.

    \subsection{Solution to The Ring}

    Let the propositional variables be:
    \[
    u = \text{the Ring has been used}, \quad
    s = \text{Sauron knows the location}, \quad
    n = \text{a Nazg\^ul knows the location}
    \]

    The initial belief state includes:

    \[
    T = Cn(\{
    \neg [g]u \wedge \neg [g]\neg u,\;
    [g](u \rightarrow s),\;
    [g](n \rightarrow s)
    \})
    \]

    \paragraph{Event.}
    Frodo uses the Ring:
    \[
    A = u
    \]

    This is a \textbf{belief update} (world change), not a revision.

    \paragraph{State-based update (KM semantics).}

    A state is a valuation over $\{u, s, n\}$.

    Initial possible worlds:
    \[
    Mod(T) = \{ s_1, s_2 \}
    \]
    where:
    \[
    s_1 = \{\neg u, \neg s, \neg n\}, \quad
    s_2 = \{\neg u, s, n\}
    \]

    We now update each state with $A = u$.

    From $s_1$:
    \[
    s_1' = \{u, \neg s, \neg n\}
    \]

    From $s_2$:
    \[
    s_2' = \{u, s, n\}
    \]

    Thus:
    \[
    Mod(T \diamond u) = \{ s_1', s_2' \}
    \]

    \paragraph{Closure under known implications.}

    From $[g](u \rightarrow s)$ and $u$, we obtain $s$ in all worlds.

    Thus:
    \[
    s_1' \rightarrow \{u, s, \neg n\}
    \]

    Final states:
    \[
    \{ \{u, s, \neg n\},\; \{u, s, n\} \}
    \]

    \paragraph{Resulting beliefs.}

    \[
    T \diamond u \vdash u \wedge s
    \]

    but:
    \[
    T \diamond u \nvdash n \quad \text{and} \quad T \diamond u \nvdash \neg n
    \]

    \paragraph{Epistemic consequences.}

    \[
    [g]u \quad \text{and} \quad [g]s
    \]

    but:
    \[
    \neg [g]n \wedge \neg [g]\neg n
    \]

    \textbf{Interpretation:}
    After the update:
    \begin{itemize}
        \item Gandalf knows that the Ring has been used.
        \item Gandalf knows that Sauron knows the location.
        \item Gandalf does not know whether a Nazg\^ul knows it.
    \end{itemize}

\end{document}